% !TeX root = ../../Book2.tex
% 纲要标题：### B.4 自旋群的平面模型与算术现象
% 纲要位置：计划纲要.md，第 1818 行。

% 纲要要点（供目录核对）：
%
% * 自旋作用的核、实平面的倍角及有理数和有限域上的像；分开分次自同构、反序对合与 Clifford 共轭，最高外次数证明行列式为一，调用中心命题求核；正式证明反射对的归一化和五元域上全部自旋像，不对一般基域点断言满射
% <!-- 短节：不设小节 -->
%

\section{自旋群的平面模型与算术现象}
\label{b2:sec:B04}

Clifford 代数的可逆元能通过共轭作用产生正交变换。这个构造还会显露底域的影响：实数域上的满射，换成有理数点后可能不再满。本节只建立所需的代数定义和一个平面模型，不预先使用 Lie 群的覆盖空间理论。仍假设 \(\operatorname{char}k\ne2\)，且 \((V,q)\) 非退化、\(\dim V\ge1\)。

\ProseParagraphBreak
把每个生成向量固定而把乘法次序颠倒，定义 Clifford 代数上的反同构
\[
\tau(v_1\cdots v_r)=v_r\cdots v_1.
\]
其良定义可由\cref{b2:thm:clifford-universal}检验：向反代数送入同一个向量仍满足 \(v^2=q(v)1\)，故唯一延伸为反同构，且平方在生成元上为恒等。\(\tau\) 称为\term[Clifford-fanxu]{反序对合}[reversion]。它固定生成向量，但对积有 \(\tau(xy)=\tau(y)\tau(x)\)。

\ProseParagraphBreak
奇偶分次给出代数自同构 \(\alpha\)：它把生成向量变为负向量，在偶部分上为恒等，在奇部分上为负恒等，且 \(\alpha(xy)=\alpha(x)\alpha(y)\)。二者的复合 \(\alpha\circ\tau\) 称为\term[Clifford-gonge]{Clifford 共轭}[Clifford conjugation]；它也把生成向量变为负向量，却反转乘法次序。\(\alpha\) 与 \(\tau\) 在生成字上交换，所以这第三种运算仍平方为恒等，并在偶部分上与 \(\tau\) 一致。许多文献先用扭转伴随作用 \(v\mapsto\alpha(g)vg^{-1}\) 定义 Clifford 群；下述自旋群只取偶元，因而 \(\alpha(g)=g\)，扭转伴随作用与普通共轭 \(v\mapsto gvg^{-1}\) 一致。

\ProseParagraphBreak
一个偶单位的共轭要成为 \(V\) 上的变换，必须把 \(V\) 保持为自身；再要求 \(\tau(g)g=1\)，则把作用元限制在反序范数为一的元素中。这两个条件分别控制共轭的作用范围和作用元的归一化。

\begin{definition}\label{b2:def:B-spin-group}
设 \(k\) 为特征不为二的域，\((V,q)\) 为维数至少一的有限维非退化 \(k\) 二次空间，\(C=\operatorname{Cl}(V,q)\)。将 \(V\) 按\cref{b2:thm:B-clifford-basis}嵌入 \(C\)，并记 \(C^0\) 为偶部分、\(\tau\) 为固定每个 \(v\in V\) 且反转乘法次序的反序对合。若偶部分中的可逆元 \(g\) 满足 \(gVg^{-1}=V\) 和 \(\tau(g)g=1\)，就把这类元素组成的集合称为\term[zixuan-qun]{自旋群}[spin group]的 \(k\) 点集，记为
\[
\operatorname{Spin}(V,q)(k)=
\bigl\{\,g\in(C^0)^\times\ \bigm|\,
gVg^{-1}=V,\ \tau(g)g=1\,\bigr\}.
\]
这里讨论的是这个具体群；若进一步把同一构造随交换 \(k\) 代数作基变换，才进入代数群的表述。
\end{definition}
\symidx{\(\operatorname{Spin}(V,q)(k)\)：自旋群的 \(k\) 点}

\begin{proposition}\label{b2:prop:B-spin-action}
设 \(k\) 为特征不为二的域，\((V,q)\) 为维数至少一的有限维非退化 \(k\) 二次空间，\(\operatorname{Spin}(V,q)(k)\) 为上一条定义的 Clifford 代数偶部分中满足正规化及反序范数为一条件的元素集合，则该集合在乘法下成群，共轭作用给出同态
\[
\rho:\operatorname{Spin}(V,q)(k)\longrightarrow\operatorname{SO}(V,q),
\qquad \rho(g)(v)=gvg^{-1},
\]
其核为 \(\{1,-1\}\)。这里不对一般底域的 \(k\) 点断言满射。
\end{proposition}
\begin{proof}
记 \(C=\operatorname{Cl}(V,q)\)。正规化 \(V\) 的条件对乘法及逆封闭。若 \(\tau(g)g=\tau(h)h=1\)，则
\[
\tau(gh)gh=\tau(h)\tau(g)gh=1,
\qquad \tau(g^{-1})g^{-1}=1.
\]
所以所给集合是群。共轭在 \(V\) 上线性且可逆；又
\[
(gvg^{-1})^2=gv^2g^{-1}=q(v)1,
\]
故它保持 \(q\)。

取正交基 \(e_1,\ldots,e_n\)，各 \(q(e_i)\ne0\)，并令 \(\omega=e_1\cdots e_n\)。把 \(e_i\) 移过其余 \(n-1\) 个正交基向量，得到 \(e_i\omega=(-1)^{n-1}\omega e_i\)。偶单项式的每个因子都贡献这个符号，而因子个数为偶数，故总符号为一；于是所有偶元素都与 \(\omega\) 交换，特别地 \(g\omega g^{-1}=\omega\)。另一方面左侧等于 \(\rho(g)(e_1)\cdots\rho(g)(e_n)\)。Clifford 乘法的关系可能产生低次数项，不能把这种乘积本身直接当作外积；进入伴随分次以后，最高次数的类才是外积。对 \(x\in F_nC\)，记 \([x]_n\) 为它在 \(F_nC/F_{n-1}C\) 中的类经\eqref{b2:eq:B-clifford-associated-graded}对应的最高外次数分量，则
\[
\begin{aligned}
[g\omega g^{-1}]_n
&=\rho(g)(e_1)\wedge\cdots\wedge\rho(g)(e_n)\\
&=\det\rho(g)\,e_1\wedge\cdots\wedge e_n,\\
[\omega]_n&=e_1\wedge\cdots\wedge e_n.
\end{aligned}
\]
第二个等号是最高外幂上的行列式展开。由于 \(g\omega g^{-1}=\omega\) 且这个外积非零，得到 \(\det\rho(g)=1\)。

由于 \(V\) 生成 \(C\)，共轭作用在 \(V\) 上为恒等，恰好意味着作用元与整个 \(C\) 交换。因此
\[
\begin{aligned}
\ker\rho
&=\operatorname{Spin}(V,q)(k)\cap Z(C)\\
&=\{\,g\in(C^0)^\times\cap Z(C)\mid\tau(g)g=1\,\}.
\end{aligned}
\]
\cref{b2:prop:B-clifford-center}给出 \(Z(C)\cap C^0=k1\)，因此核中的元素必须是非零标量。对标量 \(g\)，反序对合保持它，归一化条件便是 \(g^2=1\)。域的特征不为二，恰有两个解 \(1,-1\)，它们都满足定义中的条件，所以核为 \(\{1,-1\}\)。
\end{proof}

作为实际计算，取 \(q=\langle1,1\rangle\)，令 \(u=e_1e_2\)，则 \(u^2=-1\)、\(\tau(u)=-u\)。所有偶元素形如 \(g=a+bu\)，且
\[
\tau(g)g=a^2+b^2.
\]
若 \(a^2+b^2=1\)，则 \(g^{-1}=a-bu\)，直接相乘得到
\begin{equation}\label{b2:eq:B-spin-plane}
\rho(a+bu)=
\begin{pmatrix}a^2-b^2&2ab\\-2ab&a^2-b^2\end{pmatrix}.
\end{equation}
因此这时正规化条件自动成立，自旋群恰由 \(a^2+b^2=1\) 的偶元素组成。在 \(k=\mathbb R\) 上令 \(a=\cos\theta\)、\(b=-\sin\theta\)，则
\[
a^2-b^2=\cos2\theta,\qquad 2ab=-\sin2\theta.
\]
所以 \(g=\cos\theta-\sin\theta\,u\) 的作用是角度 \(2\theta\) 的旋转。每个平面旋转有两个互为负数的原像，倍角现象完全来自这次 Clifford 乘法计算。

\ProseParagraphBreak
在 \(k=\mathbb Q\) 上，旋转矩阵 \(\begin{psmallmatrix}0&-1\\1&0\end{psmallmatrix}\) 保持 \(\langle1,1\rangle\)，且行列式为一，所以这个有理特殊正交变换确实存在。它却没有有理自旋原像：\eqref{b2:eq:B-spin-plane}要求 \(a^2=b^2=1/2\)，而 \(1/2\) 不是有理数的平方。所缺的是满足这些条件的有理作用元，不能把实平面的满射结论无条件搬到任意域。这与\cref{b2:exm:B-rational-norm-obstruction}的范数障碍类似：实数中可解的二次方程，在有理数中可能受到平方类或同余条件的限制。

\ProseParagraphBreak
底域也决定二次方程是否有非零解。例如 \(x^2+y^2-3z^2=0\) 在实数域有解 \((\sqrt3,0,1)\)，在有理数域却只有零解：\(z=0\) 时显然 \(x=y=0\)，而 \(z\ne0\) 时除以 \(z^2\)，就得到\cref{b2:exm:B-rational-norm-obstruction}已经排除的范数方程 \((x/z)^2+(y/z)^2=3\)。这里的模三计算提供了实数方法看不到的算术障碍。

\ProseParagraphBreak
上述偶元素的共轭还可以从反射产生：一个非各向同性向量给出一次扭转共轭，两次这样的作用复合以后，两个负号相消，便成为偶元素的普通共轭。

\begin{proposition}\label{b2:prop:B-spin-reflection-pair}
设 \(k\) 为特征不为二的域，\((V,q)\) 为维数至少一的有限维非退化 \(k\) 二次空间，\(b\) 为其归一化极化型，\(C=\operatorname{Cl}(V,q)\)。若 \(z,w\in V\) 满足 \(q(z),q(w)\ne0\)，记
\[
r_z(v)=v-\frac{2b(v,z)}{q(z)}z,
\qquad r_w(v)=v-\frac{2b(v,w)}{q(w)}w
\]
为沿这两个向量的反射，则对每个 \(v\in V\) 有
\[
-zvz^{-1}=r_z(v),\qquad
(zw)v(zw)^{-1}=r_z\bigl(r_w(v)\bigr).
\]
若进一步有 \(q(z)q(w)=c^2\)、\(c\in k^\times\)，则 \(c^{-1}zw\in\operatorname{Spin}(V,q)(k)\)，且其作用为 \(r_z\circ r_w\)。
\end{proposition}
\begin{proof}
关系 \(z^2=q(z)1\) 说明 \(z^{-1}=z/q(z)\)。利用极化关系 \(zv+vz=2b(v,z)1\)，有
\[
-zvz^{-1}
=\bigl(vz-2b(v,z)1\bigr)z^{-1}
=v-\frac{2b(v,z)}{q(z)}z.
\]
\cref{b2:lem:orthogonal-reflection}说明这个公式确是反射；对 \(w\) 同样成立。因 \((zw)^{-1}=w^{-1}z^{-1}\)，将两式复合时两个负号相消，得到
\[
r_z\bigl(r_w(v)\bigr)=z w v w^{-1}z^{-1}.
\]
所以偶单位 \(zw\) 的共轭保持 \(V\)。反序对合固定 \(z,w\)，故
\[
\tau(zw)zw=wzzw=q(z)q(w)1.
\]
当这个标量为 \(c^2\) 时，\(g=c^{-1}zw\) 为偶单位，它与 \(zw\) 有相同的共轭作用，且 \(\tau(g)g=c^{-2}q(z)q(w)1=1\)。这验证了自旋群定义的两个条件。
\end{proof}

\begin{proposition}\label{b2:prop:B-finite-spin-image}
在 \(k=\mathbb F_5\) 上，令 \(V=k^2\)、\(q(xe_1+ye_2)=x^2+y^2\)，并在 \(\operatorname{Cl}(V,q)\) 中记 \(u=e_1e_2\)。则
\[
\operatorname{Spin}(V,q)(k)=\{1,-1,u,-u\},
\qquad
\rho\bigl(\operatorname{Spin}(V,q)(k)\bigr)=\{I_2,-I_2\}.
\]
相应的特殊正交群为
\[
\operatorname{SO}(V,q)
=\left\{I_2,-I_2,
\begin{pmatrix}0&-1\\1&0\end{pmatrix},
\begin{pmatrix}0&1\\-1&0\end{pmatrix}\right\}.
\]
特别地，矩阵 \(\begin{psmallmatrix}0&-1\\1&0\end{psmallmatrix}\) 不在自旋作用的像中。
\end{proposition}
\begin{proof}
由平面模型，偶元素为 \(a+bu\)，属于自旋群恰好要求 \(a^2+b^2=1\)。\(\mathbb F_5\) 中的平方是 \(0,1,4\)，其中只有 \(1+0\) 与 \(0+1\) 等于一。因此全部解为 \((a,b)=(\pm1,0),(0,\pm1)\)，给出所列四个自旋元。\eqref{b2:eq:B-spin-plane}把 \(\pm1\) 送到 \(I_2\)，把 \(\pm u\) 送到 \(-I_2\)。

若 \(A=\begin{psmallmatrix}x&y\\z&w\end{psmallmatrix}\) 保持 \(q\) 且 \(\det A=1\)，则 \(A^{\mathsf T}A=I_2\)。将 \(A^{\mathsf T}=A^{-1}\) 与行列式为一时的逆矩阵公式比较，得到 \(w=x,y=-z\)，所以
\[
A=\begin{pmatrix}x&-z\\z&x\end{pmatrix},\qquad x^2+z^2=1.
\]
同一平方计算只给出 \((x,z)=(\pm1,0),(0,\pm1)\)，正是所列四个矩阵；每个矩阵也直接满足正交及行列式为一的条件。两个非对角矩阵不等于 \(\pm I_2\)，故不在自旋作用的像中。
\end{proof}

继续研究这些问题时，自旋群的一般结构需要代数群与群作用的知识；数域二次型则要比较实、复及非 Archimedes 完备化上的解，并研究这些信息能否决定原域上的解。范数方程、平方类、Clifford 乘法及\cref{b2:prop:B-binary-norm-classification}所给的二元分类判据，已经能够处理上面的例子。关于实二次空间上的自旋群，可进一步参阅\cite[Chapter 16]{Porteous1995CliffordClassicalGroups}。
